%%%%%%%%%%%%%%%%%%%%%%%%%%%%%%%%%%%%%%%%%%%%%%%%%%%%%%%%%%%%%%%%%%%%%%
% Lớp 10 | Véctơ (chưa xét tọa độ) | Khái niệm véctơ
% Dạng 0H5?1-1: Xác định một véctơ
% Số câu: 5  (sinh tự động bởi tach_cau.py; từ nay thầy sửa trực tiếp file này)
%%%%%%%%%%%%%%%%%%%%%%%%%%%%%%%%%%%%%%%%%%%%%%%%%%%%%%%%%%%%%%%%%%%%%%

%%%%%%%%%%%%%%%% TRẮC NGHIỆM 4 PHƯƠNG ÁN %%%%%%%%%%%%%%%%

% nguồn: MaHoaNH_TracNghiem_Lop10
\begin{ex}%[0H5N1-1]%[ma:82611107]
	Cho ba điểm $A$, $B$, $C$ không thẳng hàng. Số vectơ khác vectơ không  có điểm đầu và điểm cuối trùng với các điểm $A$, $B$, $C$ là
	\choice 
		{$3$} 
		 {$5$}
		 {\True $6$}
		 {$9$}
	 \loigiai{
		 Có $6$ vectơ thoả mãn đề bài là $\overrightarrow{AB}$, $\overrightarrow{BC}$, $\overrightarrow{CA}$, $\overrightarrow{BA}$, $\overrightarrow{AC}$, $\overrightarrow{CB}$.
	 }
\end{ex}

% nguồn: MaHoaNH_TracNghiem_Lop10
\begin{ex}%[0H5N1-1]%[ma:07991500]
	Cho tứ giác  $ABCD$. Số vectơ-không nhận các đỉnh của tứ giác làm điểm đầu và điểm cuối là
	\choice 
		 {$2$}
		 {$6$}
		{\True $4$} 
		{$12$} 
	 \loigiai{
		Cứ mỗi đỉnh sẽ tạo thành một vectơ-không nên có tất cả $4$ vectơ-không nhận  các đỉnh của tứ giác làm điểm đầu và điểm cuối.
	 }
\end{ex}

% nguồn: MaHoaNH_TracNghiem_Lop10
\begin{ex}%[0H5N1-1]%[ma:1b0d86c1]
	Véc-tơ có điểm đầu là $A$, điểm cuối là $B$  kí hiệu là
	 \choice
		{\True $\overrightarrow{AB}$} 
		{$AB$} 
		 {$|\overrightarrow{AB}|$}
		 {$\overrightarrow{BA}$}
	\loigiai{ 
		Véc-tơ có điểm đầu  là $A$, điểm cuối là $B$ kí hiệu là $\overrightarrow{AB}$.
	 }
\end{ex}

% nguồn: MaHoaNH_TracNghiem_Lop10
\begin{ex}%[0H5N1-1]%[ma:ba6db726]
	Cho tam giác $ABC$. Có bao nhiêu vectơ khác vectơ-không có điểm đầu và điểm cuối là  các đỉnh $A$, $B$, $C$?
	\choice 
		{$3$} 
		{\True $6$} 
		 {$4$}
		 {$9$}
	\loigiai{ 
		Các vectơ khác vectơ-không có điểm đầu và điểm cuối là các đỉnh $A$, $B$,  $C$ là $\overrightarrow{AB}$, $\overrightarrow{BA}$, $\overrightarrow{AC}$, $\overrightarrow{CA}$, $\overrightarrow{BC}$, $\overrightarrow{CB}$.\\
		Vậy có $6$ vectơ  thỏa yêu cầu đề bài.
	} 
\end{ex}


%%%%%%%%%%%%%%%% TRẢ LỜI NGẮN %%%%%%%%%%%%%%%%

% nguồn: MaHoaNH_TraLoiNgan_Lop10
\begin{bt}%[0H5N1-1]%[ma:45029f41]
	Ta tạo được bao nhiêu vectơ khác vectơ-không có điểm đầu và điểm cuối nằm trong $4$ điểm  phân biệt cho trước trên mặt phẳng?
	\shortans{12} 
	\loigiai{ 
		Chọn điểm đầu có $4$ cách, sau đó ta chọn điểm cuối trong $3$ điểm khác  điểm đầu nên ta tạo được $3\cdot 4 = 12$ vectơ.
	 }
\end{bt}

% nguồn: TuLuan_Lop10_Dot2 (nhập 14/11/2026)
\begin{bt}%[0H5H1-1]%[ma:37628e29]
	Cho bốn điểm phân biệt không thẳng hàng $A,~B,~C,~D$. Có bao nhiêu véc-tơ khác véc-tơ không có điểm đầu và điểm cuối là hai trong bốn điểm đã cho?
	\loigiai
	{Có $4$ điểm tạo thành $6$ đoạn thẳng. Mỗi đường thẳng có $2$ véc-tơ. Vậy có $12$ véc-tơ}
\end{bt}

% nguồn: TuLuan_Lop10_Dot2 (nhập 14/11/2026)
\begin{bt}%[0H5H1-1]%[ma:56615b42]
	Cho tam giác $ABC$. Gọi $M$, $N$, $P$ lần lượt là trung điểm của $BC$, $CA$, $AB$.
	\begin{enumerate}
		\item Liệt kê tất cả các véc-tơ khác véc-tơ $\overrightarrow{0}$, cùng phương với $\overrightarrow{MN}$ và có điểm đầu, điểm cuối lấy trong các điểm đã cho.
		\item Liệt kê các véc-tơ khác véc-tơ $\overrightarrow{0}$, cùng hướng với $\overrightarrow{AB}$ và có điểm đầu, điểm cuối lấy trong các điểm đã cho.
	\end{enumerate}
	\loigiai{
		\begin{center}
			\begin{tikzpicture}[scale=0.8,>=stealth, font=\footnotesize, line join=round, line cap=round]
				\coordinate (A) at (3,5);
				\coordinate (B) at (0,0);
				\coordinate (C) at (10,0);
				\coordinate (M) at ($(C)!0.5!(B)$);
				\coordinate (N) at ($(A)!0.5!(C)$);
				\coordinate (P) at ($(A)!0.5!(B)$);
				%					\coordinate(A') at ($(A)+(-5,0)$);
				%				\coordinate(B') at ($(B)+(-5,0)$);
				\draw (A)--(B)--(C)--(A);
				\draw (M)--(N)--(P)--(M);
				\foreach 	\x/ \goc in {A/90,M/-90,C/0,B/-90,N/0,P/150} 
				\fill (\x) circle (1pt)	($(\x)+(\goc:3mm)$) node {$\x$};
			\end{tikzpicture}
		\end{center}
		\begin{enumerate}
			\item Các véc-tơ khác véc-tơ $\overrightarrow{0}$, cùng phương với $\overrightarrow{MN}$ là $\overrightarrow{NM}, \overrightarrow{AB}, \overrightarrow{BA}, \overrightarrow{AP}, \overrightarrow{PA}, \overrightarrow{BP}, \overrightarrow{PB}$.
			\item Các véc-tơ khác véc-tơ $\overrightarrow{0}$, cùng hướng với $\overrightarrow{AB}$ là $\overrightarrow{AP}, \overrightarrow{PB}, \overrightarrow{NM}$.	
		\end{enumerate}		
	}
\end{bt}

% nguồn: TuLuan_Lop10_Dot2 (nhập 14/11/2026)
\begin{bt}%[0H5N1-1]%[Trong Nguyen]%[ma:6992654d]
	Cho hình vuông $ ABCD $ có hai đường chéo cắt nhau tại $ O $. Hãy chỉ ra 
	tập hợp $ S $ gồm tất cả các vectơ khác $ \overrightarrow{0} $, có điểm đầu và điểm 
	cuối thuộc tập hợp \{$ A $; $ B $; $ C $; $ D $; $ O $\}. Hãy chia tập hợp 
	$ S $ thành các nhóm sao cho hai vectơ thuộc cùng một nhóm khi và chỉ khi 
	chúng bằng nhau.
	\loigiai{
		\immini
		{
			Ta có 
			\allowdisplaybreaks
			\begin{align}
				S =
				& \{ \overrightarrow{AB}, \overrightarrow{AC},\overrightarrow{AD}, \overrightarrow{AO}, \overrightarrow{BA},\overrightarrow{BC}, 
				\overrightarrow{BD},\overrightarrow{BO}, \overrightarrow{CA},\overrightarrow{CB}, \nonumber\\
				&\overrightarrow{CD}, \overrightarrow{CO}, \overrightarrow{DA}, \overrightarrow{DB}, \overrightarrow{DC},\overrightarrow{DO}, \overrightarrow{OA}, 
				\overrightarrow{OB}, \overrightarrow{OC}, \overrightarrow{OD} \nonumber\} 
			\end{align}
		}
		{
			\begin{tikzpicture}[>=stealth,scale=1, line join = round, line cap 
				= round]
				\coordinate (A) at (0,0);
				\coordinate (B) at (3,0);
				\coordinate (C) at (3,3);
				\coordinate (D) at ($(A)+(C)-(B)$);
				\draw (A) node[left] {$ A $}--(B)node[right] {$ B $}--
				(C)node[right] {$ C $}--(D)node[left] {$ D $}--cycle;
				\draw (A)--(C) (B)--(D) node[midway,above]{$ O $};
			\end{tikzpicture}
		}
		Các nhóm gồm các vectơ bằng nhau của tập $ S $ là 
		\begin{itemize}
			\begin{multicols}{4}
				\item $ S_1=\left\{\overrightarrow{AB}, \overrightarrow{DC} \right\}$. 
				\item $ S_2=\left\{\overrightarrow{BA}, \overrightarrow{CD} \right\}$. 
				\item $ S_3=\left\{\overrightarrow{AD}, \overrightarrow{BC} \right\}$. 
				\item $ S_4=\left\{\overrightarrow{DA}, \overrightarrow{CB} \right\}$. 
				\item $ S_5=\left\{\overrightarrow{OA}, \overrightarrow{CO} \right\}$. 
				\item $ S_6=\left\{\overrightarrow{AO}, \overrightarrow{OC} \right\}$. 
				\item $ S_7=\left\{\overrightarrow{DO}, \overrightarrow{OB} \right\}$. 
				\item $ S_8=\left\{\overrightarrow{OD}, \overrightarrow{BO} \right\}$. 
			\end{multicols}
			
		\end{itemize}
		
	}	
\end{bt}

% nguồn: TuLuan_Lop10_Dot2 (nhập 14/11/2026)
\begin{bt}%[0H5N1-1]%[Trong Nguyen]%[ma:eb654505]
	Cho đoạn thẳng $ MN $ có trung điểm là $ I $.
	\begin{enumerate}
		\item Viết các vec-tơ khác vec-tơ không có điểm đầu và điểm cuối là một 
		trong 
		ba điểm $ M $, $ N $, $ I $.
		\item Vectơ nào bằng $ \overrightarrow{MI} $ ? bằng $ \overrightarrow{NI} $ ?.
	\end{enumerate}
	\loigiai{
		\immini
		{
			\begin{enumerate}
				\item Các vec-tơ khác vec-tơ không có điểm đầu và điểm cuối là một 
				trong ba điểm $ M $, $ N $, $ I $: $ \overrightarrow{MN} $, $ \overrightarrow{NM} $, $ 
				\overrightarrow{MI} $, $ \overrightarrow{IM} $, $ \overrightarrow{NI} $, $ \overrightarrow{IN} $.
				\item 
				Các vectơ bằng $ \overrightarrow{MI} $: $ \overrightarrow{IN} $.\\
				Các vectơ bằng $ \overrightarrow{NI} $: $ \overrightarrow{IM} $.
			\end{enumerate}	
		}
		{
			\begin{tikzpicture}[>=stealth,scale=1, line join = round, line cap 
				= round]
				\coordinate (M) at (0,0);
				\coordinate (N) at (4,0);
				\coordinate (I) at ($(M)!0.5!(N)$);
				\draw (M) node[below] {$ M $}--(I)node[below] {$ I $}--
				(N)node[below] {$ N $};
				\foreach \diem in {M,N,I}
				\draw[fill=black] (\diem) circle(1pt);
			\end{tikzpicture}
		}
		
	}	
\end{bt}

% nguồn: TuLuan_Lop10_Dot2 (nhập 14/11/2026)
\begin{bt}%[0H5H1-1]%[ma:2b16993e]
		Cho hình bình hành $ABCD$. Hãy chỉ ra các véc-tơ (khác $\overrightarrow{0}$) có điểm đầu và điểm cuối là một trong bốn đỉnh của hình bình hành $ABCD$. Trong số các véc-tơ trên, hãy chỉ ra
		\begin{enumerate}
			\item Các véc-tơ cùng phương.
			\item Các cặp véc-tơ cùng phương nhưng ngược hướng.
			\item Các cặp véc-tơ bằng nhau.
		\end{enumerate}
		\loigiai
		{
			\immini
			{
				Các véc-tơ (khác $\overrightarrow{0}$) có điểm đầu và điểm cuối là một trong bốn đỉnh của hình bình hành $ABCD$ là $\overrightarrow{AB}$, $\overrightarrow{BA}$, $\overrightarrow{AC}$, $\overrightarrow{CA}$, $\overrightarrow{AD}$, $\overrightarrow{DA}$, $\overrightarrow{BC}$, $\overrightarrow{CB}$, $\overrightarrow{BD}$, $\overrightarrow{DB}$, $\overrightarrow{CD}$, $\overrightarrow{DC}$.
			}
			{
				\begin{tikzpicture}
					\tkzDefPoints{0/0/A, 3/0/B, 4/1/C}
					\tkzDefPointBy[translation=from B to A](C)\tkzGetPoint{D}
					\tkzDrawPoints[fill=black](A,B,D,C)
					\tkzDrawPolygon(A,B,C,D)
					\tkzLabelPoints[left](A,D)
					\tkzLabelPoints[right](B,C)
				\end{tikzpicture}
			}
			\noindent
			\begin{enumerate}
				\item Các véc-tơ cùng phương với nhau là
				\begin{itemize}
					\item $\overrightarrow{AB}$, $\overrightarrow{BA}$, $\overrightarrow{CD}$, $\overrightarrow{DC}$.
					\item $\overrightarrow{AD}$, $\overrightarrow{DA}$, $\overrightarrow{BC}$, $\overrightarrow{CB}$.
					\item $\overrightarrow{AC}$ và $\overrightarrow{CA}$.
					\item $\overrightarrow{BD}$ và $\overrightarrow{DB}$.
				\end{itemize}
				\item Các cặp véc-tơ cùng phương nhưng ngược hướng là $\overrightarrow{AB}$ và $\overrightarrow{BA}$; $\overrightarrow{BC}$ và $\overrightarrow{CB}$; $\overrightarrow{CD}$ và $\overrightarrow{DC}$; $\overrightarrow{DA}$ và $\overrightarrow{AD}$; $\overrightarrow{AC}$ và $\overrightarrow{CA}$; $\overrightarrow{BD}$ và $\overrightarrow{DB}$; $\overrightarrow{AB}$ và $\overrightarrow{CD}$; $\overrightarrow{BA}$ và $\overrightarrow{DC}$; $\overrightarrow{AD}$ và $\overrightarrow{CB}$; $\overrightarrow{DA}$ và $\overrightarrow{BC}$.
				\item Các cặp véc-tơ bằng nhau là $\overrightarrow{AB}$ và $\overrightarrow{DC}$; $\overrightarrow{BA}$ và $\overrightarrow{CD}$; $\overrightarrow{AD}$ và $\overrightarrow{BC}$; $\overrightarrow{DA}$ và $\overrightarrow{CB}$.
			\end{enumerate}
		}
	\end{bt}

% nguồn: TuLuan_Lop10_Dot2 (nhập 14/11/2026)
\begin{bt}%[0H5N1-1]%[ma:e298ca30]
	Cho hai điểm phân biệt $A,~B$. Tìm số véc-tơ được tạo thành từ hai điểm này?
	\loigiai{
	
	Véc-tơ là đoạn thẳng được định hướng. Như vậy từ hai điểm $A,B$ ta được hai véc-tơ là $\overrightarrow{AB}$ và $\overrightarrow{BA}$.
}
\end{bt}

% nguồn: TuLuan_Lop10_Dot2 (nhập 14/11/2026)
\begin{bt}%[0H5N1-1]%[ma:ce8c9fc6]
	Cho tam giác $ABC$. Viết tất cả các véc-tơ mà điểm đầu và điểm cuối là $A$, $B$ hoặc $C$.
	\loigiai{
		Ta có $\overrightarrow{AB}$, $\overrightarrow{BA}$, $\overrightarrow{AC}$, $\overrightarrow{CA}$, $\overrightarrow{CB}$, $\overrightarrow{BC}$.
	}
	\end{bt}

% nguồn: TuLuan_Lop10_Dot2 (nhập 14/11/2026)
\begin{bt}%[0H5H1-1]%[Dự án đề kiểm tra HKII NH22-23-BCTuan]%[THPT Thực Hành Sài Gòn]%[ma:c9233228]
	\immini
	{
		Cho $6$ điểm cùng nằm trên một đường tròn như hình bên.
		\begin{enumerate}
			\item Có bao nhiêu véc-tơ có điểm đầu và điểm cuối thuộc các điểm đã cho?
			\item Có bao nhiêu tam giác có đỉnh thuộc các điểm đã cho?
		\end{enumerate}
	}
	{
	\begin{tikzpicture}[line join=round, line cap=round,>=stealth,font=\footnotesize,scale=0.5]
\def\r{4}
\coordinate (O) at (0,0);
\draw (O) circle (\r);
\foreach \i/\j/\k in {0/E/0,1/F/90,2/A/90,3/B/180,4/C/-90,5/D/-90}
\draw ($(O)+(60*\i:4)$) circle (1pt) +(\k:0.5cm) node {$\j$};
\end{tikzpicture}
	}
\loigiai{
\begin{enumerate}
	\item Số véc-tơ có điểm đầu và điểm cuối thuộc các điểm đã cho là $\mathrm{A}_6^2=30$.
	\item Số tam giác có các đỉnh thuộc các điểm đã cho là $\mathrm{C}_6^3=20$.
\end{enumerate}
}
\end{bt}

% nguồn: TuLuan_Lop10_Dot2 (nhập 14/11/2026)
\begin{bt}%[0H5H1-1]%[ma:5ce53114]
		Cho hai điểm phân biệt $H$, $K$. Viết hai véc-tơ mà điểm đầu và điểm cuối là $H$ hoặc $K$.
	\loigiai{
		Hai véc-tơ thỏa mãn yêu cầu đề bài là $\overrightarrow{HK}$ và $\overrightarrow{KH}$.
	}
\end{bt}

% nguồn: TuLuan_Lop10_Dot2 (nhập 14/11/2026)
\begin{bt}%[0H5N1-1]%[Dự án 4 - TLDH - TeamTeXHoa - Lê Quân]%[ma:f6c505d9]
Cho tứ giác $ABCD$. Có bao nhiêu véc-tơ khác véc-tơ-không có điểm đầu và điểm cuối là đỉnh của tứ giác.
	\loigiai
	{
	Hai điểm phân biệt, chẳng hạn $A,\,\,B$ ta xác định được hai véc-tơ khác véc-tơ-không là $\overrightarrow{AB},\,\,\overrightarrow{BA}$.\\
	 Mà từ bốn đỉnh $A,\,\,B,\,\,C,\,\,D$ của tứ giác ta có 6 cặp điểm phân biệt do đó có 12 véc-tơ thỏa mãn yêu cầu bài toán.	
	}
\end{bt}

% nguồn: TuLuan_Lop10_Dot2 (nhập 14/11/2026)
\begin{bt}%[0H5N1-1]%[Dự án bài tập sách CTST]%[Phan Anh]%[ma:6c514302]
	\immini{Cho hình vuông $ABCD$ có tâm $O$ và có cạnh bằng $1$ (tham khảo hình vẽ).
		\begin{enumerate}
			\item Tìm điểm đầu, điểm cuối, giá và độ dài của các véc-tơ: $\overrightarrow{CA}$, $\overrightarrow{OA}$, $\overrightarrow{BD}$.
			\item Tìm các véc-tơ đơn vị trong hình.
	\end{enumerate}}
	{\begin{tikzpicture}[>=stealth,line join=round,line cap=round,font=\footnotesize,scale=0.7]
			\path (0,3) coordinate (A)
			(3,3) coordinate (B)
			(3,0) coordinate (C)
			(0,0) coordinate (D)
			($(A)!0.5!(C)$)coordinate (O)
			;
			\draw (C)--(A)--(B)--(C)--(D)--(A) (B)--(D);
			\foreach \p / \r in {A/135,B/45,C/-45,D/-135,O/-90}
			\fill (\p) circle (1.2pt) node[shift={(\r:3mm)}]{$\p$};
			\tkzLabelSegments[left,pos=0.5](A,D){$1$}
	\end{tikzpicture}}
	\loigiai{
		\begin{enumerate}
			\item Véc-tơ $\overrightarrow{CA}$ có điểm đầu là $C$, điểm cuối là $A$ và có giá là đường thẳng $AC$.\\
			Véc-tơ $\overrightarrow{OA}$ có điểm đầu là $O$, điểm cuối là $A$ và có giá là đường thẳng $AO$.\\
			Véc-tơ $\overline{BD}$ có điểm đầu là $B$, điểm cuối là $D$ và có giá là đường thẳng $BD$.\\
			Ta có: $CA=BD=\sqrt{2};OA=\dfrac{\sqrt{2}}{2}$.\\
			Suy ra $|\overrightarrow{CA}|=|\overrightarrow{BD}|=\sqrt{2};|\overrightarrow{OA}|=\dfrac{\sqrt{2}}{2}$.
			\item Các véc-tơ đơn vị là $\overrightarrow{AB}$, $\overrightarrow{BA}$, $\overrightarrow{BC}$, $\overrightarrow{CB}$, $\overrightarrow{CD}$, $\overrightarrow{DC}$, $\overrightarrow{DA}$, $\overrightarrow{AD}$.
	\end{enumerate}}
\end{bt}

% nguồn: TuLuan_Lop10_Dot2 (nhập 14/11/2026)
\begin{bt}%[0H5N1-1]%[Dự án bài tập sách CTST]%[Phan Anh]%[ma:925ef1af]
	Bạn hãy tìm sự khác biệt giữa hai đại lượng sau
	\begin{itemize}
		\item Chiếc xe máy có giá tiền là $30$ triệu đồng.
		\item Chiếc thuyền chạy với vận tốc là $30$km/h theo hướng tây nam.
	\end{itemize}
	\loigiai{
		Một đại lượng chi số lượng; đại lượng còn lại chỉ rõ giá trị và hướng.}
\end{bt}

% nguồn: TuLuan_Lop10_Dot2 (nhập 14/11/2026)
\begin{bt}%[0H5N1-1]%[Dự án bài tập sách CTST]%[Phan Anh]%[ma:0f19612e]
	Trong các đại lượng sau: nhiệt độ, lực, thể tích, tuổi, độ dịch chuyển, vận tốc, đại lượng nào cần được biểu diễn bởi véc-tơ? 
	\loigiai{
		Các đại lượng cần được biểu diễn bởi véc-tơ là lực, sức hút, độ dịch chuyển, vận tốc.}
\end{bt}

% nguồn: chuyen_dang (nhập 22/11/2026)
\begin{ex}%[0H5N1-1]%[ma:2c0bbf3d]
	Cho hình lục giác đều $ABCDEF$ tâm $O$. Số các vectơ khác vectơ không, cùng phương với vectơ $\overrightarrow{OB}$ có điểm đầu và điểm cuối là các đỉnh của lục giác  là
	\choice 
		{$4$} 
		{\True  $6$}
		{$8$} 
		{$10$} 
	 \loigiai{
		\begin{center}
			\begin{tikzpicture}[scale=1,line join=round,line cap=round,font=\footnotesize,>=stealth]
					\def\a{2} 
					\path (0:0) coordinate  (O)
						(0:\a) coordinate (C) 
						(60:\a) coordinate (D) 
						(120:\a)  coordinate (E)
						(180:\a)  coordinate (F)
						(240:\a)  coordinate (A)
						(300:\a)  coordinate (B);
					\draw  (A)--(B)--(C)--(D)--(E)--(F)--(A)
						(O)--(B); 
					\foreach \x/  \goc in {A/180,B/-135,C/-45,D/0,E/90,F/180,O/0}
					\fill (\x) circle (1pt) 
						 ($(\x)+(\goc:3mm)$) node {$\x$};
				\end{tikzpicture}
		\end{center}
		Các vectơ cùng  phương với vectơ $\overrightarrow{OB}$ là\\
		$\overrightarrow{BE}$, $\overrightarrow{EB}$,  $\overrightarrow{DC}$, $\overrightarrow{CD}$, $\overrightarrow{FA}$, $\overrightarrow{AF}$.
	} 
\end{ex}

% nguồn: chuyen_dang (nhập 22/11/2026)
\begin{bt}%[0H5N1-1]%[Dự án 4 - TLDH - TeamTeXHoa - Lê Quân]%[ma:4be064ae]
Cho tam giác $ABC$. Gọi $M,\,N,\,P$ lần lượt là trung điểm của $BC,\,CA,\,AB$.
\begin{enumerate}
\item Xác định các véc-tơ khác véc-tơ-không cùng phương với $\overrightarrow{MN}$ có điểm đầu và điểm cuối lấy trong các điểm đã cho.
\item Xác định các véc-tơ khác véc-tơ-không cùng hướng với $\overrightarrow{AB}$ có điểm đầu và điểm cuối lấy trong điểm đã cho.
\item Vẽ các véc-tơ bằng véc-tơ $\overrightarrow{NP}$ mà có điểm đầu $A,\,B$.
\end{enumerate}
\loigiai{
\immini{
\begin{enumerate}
\item Các véc-tơ khác véc-tơ không cùng phương với $\overrightarrow{MN}$ là $\overrightarrow{NM},\,\overrightarrow{AB},\,\overrightarrow{BA},\,\overrightarrow{AP},\,\overrightarrow{PA},\,\overrightarrow{BP},\,\overrightarrow{PB}$.
\item Các véc-tơ khác véc-tơ-không cùng hướng với $\overrightarrow{AB}$ là $\overrightarrow{AP},\,\,\overrightarrow{PB},\,\,\overrightarrow{NM}$ .
\item Trên tia $CB$ lấy điểm $B'$ sao cho $BB'=NP$.\\
Khi đó ta có $\overrightarrow{BB'}$ là véc-tơ có điểm đầu là $B$ và bằng véc-tơ $\overrightarrow{NP}$.\\
Qua $A$ dựng đường thẳng song song với đường thẳng $NP$. Trên đường thẳng đó lấy điểm $A'$ sao cho $\overrightarrow{AA'}$ cùng hướng với $\overrightarrow{NP}$ và $AA'=NP$.\\
Khi đó ta có $\overrightarrow{AA'}$ là véc-tơ có điểm đầu là $A$ và bằng véc-tơ $\overrightarrow{NP}$.
\end{enumerate}}
{\begin{tikzpicture}[line join=round, line cap=round,thick, scale=.7]
	\coordinate (A) at (0,3);
	\coordinate (B) at (-2,0);
	\coordinate (C) at (3,0);
	\draw(A)--(B)--(C)--cycle;
	\coordinate (P) at ($(A)!0.5!(B)$);
	\coordinate (N) at ($(A)!0.5!(C)$);
	\coordinate (M) at ($(C)!0.5!(B)$);
	\coordinate (B') at ($(B)+(P)-(N)$);
		\coordinate (A') at ($(A)+(P)-(N)$);
		\draw(M)--(N)--(P)--cycle;
		\draw[->] (B)--(B');
		\draw[->] (A)--(A');
	\foreach \i/\g in {A/90, A'/90, B'/90, M/-90, P/180, B/-90,C/-90, N/30}{\draw[fill=black](\i) circle (1.5pt) ($(\i)+(\g:5mm)$) node[scale=1]{$\i$};}
\end{tikzpicture}}

}
\end{bt}

% nguồn: chuyen_dang (nhập 22/11/2026)
\begin{ex}%[0H5N1-1]%[ma:b4a1849d]
	Trên mặt phẳng cho lục giác $ABCDEF$. Hỏi có  bao nhiêu véc-tơ khác véc-tơ $\overrightarrow 0 $ và có điểm đầu, điểm cuối là 2 trong 6 đỉnh của lục giác đã cho?
	\choice 
		{$36$} 
		 {\True $30$}
		{$11$} 
		{$25$} 
	 \loigiai{
	Mỗi véc-tơ tạo thành là một chỉnh hợp chập $2$ của $6$ phần tử.\\ 
	Vậy số véc-tơ thoả mãn bài toán là $\mathrm{A}_6^2=  30$.
	} 
\end{ex}

% nguồn: chuyen_dang (nhập 22/11/2026)
\begin{ex}%[0H5H1-1]%[ma:bce5672c]
	Cho $n$ điểm phân biệt. Xét tất cả các véc-tơ  khác $\overrightarrow{0}$, có điểm đầu và điểm cuối là các điểm đã cho. Số véc-tơ thoả mãn là
	 \choice
		 {$n$}
		{$n-1$} 
		{$\dfrac{n(n-1)}{2}$} 
		 {\True $n(n-1)$}
	\loigiai{ 
		Mỗi véc-tơ khác $\overrightarrow{0}$  là một chỉnh hợp chập $2$ của $n$ điểm đã cho.\\
		Vậy số véc-tơ thoả  mãn bài toán là $\mathrm{A}_n^2=n(n-1)$.
	} 
\end{ex}

% nguồn: chuyen_dang (nhập 22/11/2026)
\begin{ex}%[0H5N1-1]%[ma:8f02be3c]
	 Trong mặt phẳng cho $6$ điểm phân biệt. Có bao nhiêu véc-tơ khác $\overrightarrow{0}$ có điểm đầu và điểm cuối là $2$ trong $6$ điểm đã cho?
	 \choice
		 {$1440$}
		 {\True $30$}
		{$12$} 
		 {$15$}
	\loigiai{Số véc-tơ khác  $\overrightarrow{0}$ có điểm đầu và điểm cuối là $2$ trong $6$ điểm đã cho là $\mathrm{A}_6^2=30$.
	 }
\end{ex}

% nguồn: chuyen_dang (nhập 22/11/2026)
\begin{ex}%[0H5N1-1]%[ma:9ce00d02]
	Trong mặt phẳng cho tập hợp gồm $25$ điểm phân biệt. Có bao nhiêu véc-tơ khác véc-tơ $\overrightarrow{0}$ có điểm đầu và điểm cuối thuộc  tập hợp điểm này?
	\choice 
		 {$50$}
		 {$300$}
		{\True $600$} 
		{$625$} 
	 \loigiai{
		Số véc-tơ khác $\overrightarrow{0}$ có điểm đầu và điểm cuối được tạo ra bởi $25$ điểm phân  biệt là $\mathrm{A}_{25}^2=600$ véc-tơ.
	 }
\end{ex}

% nguồn: chuyen_dang (nhập 22/11/2026)
\begin{bt}%[0H5H1-1]%[Đề kiểm tra GK1 môn Toán 10 trường THPT Đoàn Thượng]%[Mui Doan, 10EX-GK1-2223]%[ma:dac69202]
Cho hình bình hành $ABCD$ tâm $O$. Gọi $P$, $Q$, $R$ lần lượt là trung điểm của $AB$, $BC$, $AD$. Lấy $8$ điểm trên là điểm đầu hoặc điểm cuối của các véc-tơ. Tìm số các véc-tơ khác véc-tơ không, bằng với véc-tơ $\overrightarrow{AR}$.
	\loigiai{
	\immini{
	Số các véc-tơ khác véc-tơ không, bằng với véc-tơ $\overrightarrow{AR}$ là $\overrightarrow{RD}$, $\overrightarrow{BQ}$, $\overrightarrow{QC}$, $\overrightarrow{PO}$.
	}	
{
\begin{tikzpicture}[font=\footnotesize, line join=round, line cap=round, >=stealth]
\path (0,0) coordinate (A)
(-2,-2) coordinate (B)
(5,0) coordinate (D)
($(B)+(D)-(A)$) coordinate (C)
($(A)!.5!(C)$) coordinate (O)
($(A)!.5!(B)$) coordinate (P)
($(B)!.5!(C)$) coordinate (Q)
($(A)!.5!(D)$) coordinate (R)
;

\draw (A)--(B)--(C)--(D)--cycle (A)--(C) (B)--(D)
;	
\foreach \t/\g in {A/120,B/180,C/0,D/0,O/90,P/180,Q/-90,R/90}{
	\draw[fill=white] (\t) circle (1pt) node[shift={(\g:7pt)},font=\scriptsize]{$ \t $};
}
\end{tikzpicture}
	}
	}
\end{bt}

% nguồn: chuyen_dang (nhập 22/11/2026)
\begin{ex}%[0H5N1-1]%[ma:42cbd57e]
	Cho lục giác $ABCDEF$. Chọn khẳng định đúng trong các khẳng  định sau.
	\choiceTF 
		{\True Số véc-tơ khác $\overrightarrow{0}$ có điểm đầu, điểm cuối là hai  trong $6$ đỉnh của lục giác là $30$}
		{Số véc-tơ khác  $\overrightarrow{0}$ có điểm đầu, điểm cuối là hai trong $6$ đỉnh của lục giác là $15$}
		{\True Số tam  giác tạo thành từ $6$ đỉnh của lục giác là $20$}
		{Số tam giác tạo thành từ $6$ đỉnh của  lục giác là $120$}
	 \loigiai{
		\begin{itemize}
			\item Số véc-tơ khác $\overrightarrow{0}$ có điểm đầu, điểm cuối là hai trong $6$ đỉnh của  lục giác là $\mathrm{A}_6^2=30$.
			\item Số tam giác tạo thành từ  $6$ đỉnh của lục giác là $\mathrm{C}_6^3=20$.
		\end{itemize}
	} 
\end{ex}

% nguồn: chuyen_dang (nhập 22/11/2026)
\begin{ex}%[0H5H1-1]%[ma:68ae002a]
	Trong  mặt phẳng cho một tập hợp gồm $6$ điểm phân biệt. Có bao nhiêu vectơ khác vectơ $\overrightarrow 0$ có điểm đầu và điểm cuối thuộc tập hợp điểm này?
	\choice 
		 {$15$}
		{$12$} 
		 {$1\,440$}
		 {\True $30$}
	\loigiai{ 
		Mỗi cặp sắp thứ tự gồm hai điểm $\left(A, B\right)$ cho ta một vectơ có điểm đầu $A$ và điểm cuối $B$ và ngược lại. Như vậy, mỗi vectơ có thể xem là một chỉnh hợp chập $2$ của tập hợp  $6$ điểm đã cho.\\
		Suy ra có $\mathrm{A}_6^2=30$  cách.
	} 
\end{ex}

% nguồn: chuyen_dang (nhập 22/11/2026)
\begin{bt}%[0H5H1-1]%[ma:8730da09]
	Trong mặt phẳng cho một  tập hợp gồm $6$ điểm phân biệt. Có bao nhiêu vectơ khác vectơ $\overrightarrow{0}$ có điểm đầu và điểm cuối thuộc tập hợp điểm này?
	\loigiai{ 
		Mỗi vectơ thỏa đề là một chỉnh hợp chập $2$ của $6$ nên  có $\mathrm{A}_6^2 = 30$ vectơ.
	 }
\end{bt}

% nguồn: chuyen_dang (nhập 22/11/2026)
\begin{bt}%[0H5N1-1]%[ma:1a159681]
% nguon: Archive/2024/0. [LVH] BG - 12 New - 4in1-20240625T224301Z-001/0. [LVH] BG - 12 New - 4in1/data/data_C2/Bai-1-Dang-1-Cauhoilythuyet.tex #1724
[Cùng Khám phá, Mức độ 1]%[BG - 12 New - 4in1,Lam Huy]
	Cho hình chóp tứ giác $S.ABCD$. Hãy chỉ ra tất cả những véc-tơ có điểm đầu và điểm cuối là hai điểm phân biệt lấy trong các điểm $S$, $A$, $B$, $C$, $D$.
	\loigiai{
	\immini{	Ta có $20$ véc-tơ $\overrightarrow{SA}$, $\overrightarrow{AS}$, $\overrightarrow{SB}$, $\overrightarrow{BS}$, $\overrightarrow{SC}$, $\overrightarrow{CS}$, $\overrightarrow{SD}$, $\overrightarrow{DS}$, $\overrightarrow{AB}$, $\overrightarrow{BA}$, $\overrightarrow{AC}$, $\overrightarrow{CA}$, $\overrightarrow{AD}$, $\overrightarrow{DA}$, $\overrightarrow{BC}$, $\overrightarrow{CB}$, $\overrightarrow{BD}$, $\overrightarrow{DB}$, $\overrightarrow{CD}$, $\overrightarrow{DC}$.}{
		\begin{tikzpicture}[line join=round, line cap=round,thick, >=stealth, scale=0.7]
			\coordinate (S) at (1,4);
			\coordinate (A) at (0,0);
			\coordinate (B) at (1,-2);
			\coordinate (C) at (3,-2.7);
			\coordinate (D) at (5,0);
			\draw(S)--(A) (S)--(B) (S)--(C) (S)--(D) (B)--(C) (A)--(B)--(C)--(D);
			\draw[dashed,thin](A)--(D);
			\foreach \i/\g in {S/90,A/180,B/-110,C/-90,D/0}{\draw[fill=black](\i) circle (1.5pt) ($(\i)+(\g:3mm)$) node[scale=1]{$\i$};}
	\end{tikzpicture}}
	}
 \end{bt}

% nguồn: chuyen_dang (nhập 22/11/2026)
\begin{bt}%[0H5N1-1]%[ma:c8c7ee4f]
% nguon: Archive/2024/0. [LVH] BG - 12 New - 4in1-20240625T224301Z-001/0. [LVH] BG - 12 New - 4in1/data/data_C2/Bai-1-Dang-1-Cauhoilythuyet.tex #1738
[Mức độ 1]%[BG - 12 New - 4in1,Lam Huy]
	Cho tứ diện $ABCD$.
 \begin{listEX}[1]
        \item  Hãy chỉ ra các vectơ có điểm đầu là $B$ và điểm cuối là một trong các đỉnh còn lại của tứ diện.
        \item  Hãy chỉ ra các vectơ có điểm cuối là $D$ và điểm đầu là một trong các đỉnh còn lại của tứ diện.
        \item  Hãy chỉ ra các vectơ khác vectơ-không có điểm đầu và điểm cuối là một trong các đỉnh của tứ diện.
 \end{listEX}
	\loigiai{
	\immini{\begin{listEX}[1]
 \item Ngoài đỉnh $B$, tứ diện còn có $3$ đỉnh $A$, $C$, $D$ nên ta có $3$ vectơ $\overrightarrow{BA}$, $\overrightarrow{BC}$, $\overrightarrow{BD}$.
 \item Ngoài đỉnh $D$, tứ diện còn có $3$ đỉnh $A$, $B$, $C$ nên ta có $3$ vectơ $\overrightarrow{AD}$, $\overrightarrow{BD}$, $\overrightarrow{CD}$.
 \item Các vectơ có điểm đầu $A$ là $\overrightarrow{AB}$, $\overrightarrow{AC}$, $\overrightarrow{AD}$.\\
 Các vectơ có điểm đầu $B$ là $\overrightarrow{BA}$, $\overrightarrow{BC}$, $\overrightarrow{BD}$.\\
 Các vectơ có điểm đầu $C$ là $\overrightarrow{CA}$, $\overrightarrow{CB}$, $\overrightarrow{CD}$.\\
 Các vectơ có điểm đầu $D$ là $\overrightarrow{DA}$, $\overrightarrow{DB}$, $\overrightarrow{DC}$.
 \end{listEX}}{
		\begin{tikzpicture}[line join=round, line cap=round,thick, >=stealth, scale=0.7]
				\coordinate (B) at (-3,0);
				\coordinate (C) at (-1,-2);
				\coordinate (D) at (3,0);
				\coordinate (M) at ($(C)!0.5!(B)$);
				\coordinate (O) at ($(D)!2/3!(M)$);
				\coordinate (A) at ($(O)+(0,5)$);
				\foreach \x/\y in {A/B, A/D, A/C}
				\draw (\x)--(\y);
				\draw[dashed,thin](D)--(B);
				\draw (B)--(C)--(D);
				\foreach \i/\g in {A/90, D/0, B/180, C/-90}
				\draw ($(\i)+(\g:3mm)$) node[scale=1]{$\i$};
		\end{tikzpicture}}
	}
 \end{bt}

% nguồn: chuyen_dang (nhập 22/11/2026)
\begin{bt}%[0H5N1-1]%[ma:82715fa7]
% nguon: Archive/2024/0. [LVH] BG - 12 New - 4in1-20240625T224301Z-001/0. [LVH] BG - 12 New - 4in1/data/data_C2/Bai-1-Dang-1-Cauhoilythuyet.tex #1723
[Cùng Khám phá, Mức độ 1]%[BG - 12 New - 4in1,Lam Huy]
	Cho tứ diện $ABCD$. Hãy chỉ ra các véc-tơ có điểm đầu là $A$ và điểm cuối là một trong các đỉnh còn lại của tứ diện.
	\loigiai{
	\immini{Ngoài đỉnh $A$, tứ diện còn có $3$ đỉnh $B$, $C$, $D$ nên ta có $3$ véc-tơ $\overrightarrow{AB}$, $\overrightarrow{AC}$, $\overrightarrow{AD}$.}{
		\begin{tikzpicture}[line join=round, line cap=round,thick, >=stealth, scale=0.7]
				\coordinate (B) at (-3,0);
				\coordinate (C) at (-1,-2);
				\coordinate (D) at (3,0);
				\coordinate (M) at ($(C)!0.5!(B)$);
				\coordinate (O) at ($(D)!2/3!(M)$);
				\coordinate (A) at ($(O)+(0,5)$);
				\foreach \x/\y in {A/B, A/D, A/C}
				\draw[->, very thick] (\x)--(\y);
				\draw[dashed,thin](D)--(B);
				\draw (B)--(C)--(D);
				\foreach \i/\g in {A/90, D/0, B/180, C/-90}
				\draw ($(\i)+(\g:3mm)$) node[scale=1]{$\i$};
		\end{tikzpicture}}
	}
 \end{bt}
